         \chapter{Getal patrone}
    \setcounter{figure}{1}
    \setcounter{subfigure}{1}
    \label{d2d281c6991a4da87e4a19596c1ef625}
$ \hspace{-5pt}\begin{array}{cccccccccccc}   \includegraphics[width=0.75cm]{col11306.imgs/summary_fullmarks.png} &   \includegraphics[width=0.75cm]{col11306.imgs/summary_video.png} &   \end{array} $ \hspace{2 pt}\raisebox{-5 pt}{} {(section shortcode: MG10022 )} \par 

%             
In vorige jare het jy patrone gesien in die vorm van prentjies en getalle. In hierde hoofstuk sal ons meer leer
van die wiskunde van patrone. Patrone is herkenbaar as herhalende reekse wat gevind kan word in die natuur,
vorms, gebeure, groepe van getalle en op baie ander plekke in ons daaglikse lewe. Byvoorbeeld, patrone kan
gevind word in die sade van sonneblomme, sneeuvlokkies, geometriese patrone op lappieskomberse en teëls en
reekse getalle soos 
\begin{enumerate}[noitemsep, label=\textbf{\arabic*}. ] 
    \item $2;~4;~6;~8;~10;\ldots$
    \item $1;~2;~4;~7;~11;\ldots$
    \item $1;~4;~9;~16;~25;\ldots$
    \item $5;~10;~20;~40;~80;\ldots$
\end{enumerate}

   
\section{Byvoorbeeld van getal patrone}
Reekse van getalle kan interessante patrone bevat. Die volgende is ’n lys van die mees algemene patrone en
hoe hulle gevorm word.
Voorbeelde:
\begin{enumerate}[noitemsep, label=\textbf{\arabic*}. ] 
    \item $1;~4;~7;~10;~13;~16;~19;~22;~25;~\ldots$
    Hierdie reeks het ’n verskil van $3$ tussen al die getalle. Die patroon word gevorm deur elke keer.\\
     $3$ by te tel by die vorige getal
    \item $3;~8;~13;~18;~23;~28;~33;~38;~\ldots$
     Hierdie reeks het ’n verskil van $5$ tussen al die getalle. Die patroon word gevorm.\\
    deur elke keer $5$ by te tel by die vorige getal.
    \item $2;~4;~8;~16;~32;~64;~128;~256;\ldots$
    Hierdie reeks het ’n faktor van $2$ tussen al die getalle. Die volgende getal in.\\
    die reeks word gevorm deur die vorige een met $2$te vermenigvuldig.
    \item $3;~9;~27;~81;~243;~729;~2187;\ldots$
    Hierdie reeks het ’n faktor van $3$ tussen al die getalle. Die volgende getal in.\\
    die reeks word gevorm deur die vorige een met  $3$te vermenigvuldig.
\end{enumerate}
     

\subsection*{Spesiale Reekse}

\subsubsection*{Driehoeksgetalle}

% \psset{xunit=1.0cm,yunit=1.0cm,algebraic=true,dotstyle=o,dotsize=3pt 0,linewidth=0.8pt,arrowsize=3pt 2,arrowinset=0.25}
\begin{pspicture*}(-5.3,-3.74)(7.32,6.48)
\rput[tl](-5.1,0.54){$1$}
\rput[tl](-2.6,0.54){$3$}
\rput[tl](0.94,0.5){$6$}
\rput[tl](5.42,0.54){$10$}
\psdots[dotstyle=*](0,1)
\psdots[dotstyle=*](1,1)
\psdots[dotstyle=*](2,1)
\psdots[dotstyle=*](1.52,2)
\psdots[dotstyle=*](0.6,2)
\psdots[dotstyle=*](1.14,2.98)
\psdots[dotstyle=*](-5,1)
\psdots[dotstyle=*](-3,1)
\psdots[dotstyle=*](-2,1)
\psdots[dotstyle=*](-2.46,2)
\psdots[dotstyle=*](7,1)
\psdots[dotstyle=*](4,1)
\psdots[dotstyle=*](5,1)
\psdots[dotstyle=*](6,1)
\psdots[dotstyle=*](5.46,2)
\psdots[dotstyle=*](4.48,2)
\psdots[dotstyle=*](6.32,2)
\psdots[dotstyle=*](6,3)
\psdots[dotstyle=*](5,3)
\psdots[dotstyle=*](5.54,3.96)
\end{pspicture*}
$\ldots$\\
\\
$1;~3;~6;~10;~15;~21;~28;~36;~45;\ldots$\\
\\
Hierdie reekse word gevorm deur ’n patroon van kolletjies wat ’n driehoek vorm. Deur nog ’n ry kolletjies aan te
heg (waar die elke nuwe ry een meer kolletjie bevat as die vorige een) en die kolletjies te tel, is dit moontlik om
die volgende getal in die reeks te vind.\par 

\subsubsection*{Vierkantsgetalle}
$1;~4;~9;~16;~25;~36;~49;~64;~81;\ldots$\\
Die waarde van ’n term in die reeks word gevind deur die posisie (pleknommer in die ry) te kwadreer.\\
Die $2^{nd}$ getal is ${2}^{2} = 4$.
Die $7^{th}$ getal is ${7}^{2} = 49$. Wat sal die $10^{th}$ getal wees?
            
\subsubsection*{Derdemagsgetalle}
$1;~8;~27;~64;~125;~216;~343;~512;~729;\ldots$\\
Die waarde van ’n term in die reeks word gevind deur die posisie tot die derde mag te verhef.\\
Die $2^{nd}$ getal is $2^{3}=8$.
Die $7^{th}$ getal is $7^{3}=343$. Wat sal die $10^{th}$ getal wees?

\subsubsection*{Fibonacci Getalle}
$0;~1;~1;~2;~3;~5;~8;~13;~21;~34;...$\\
Die waarde van ’n term in die reeks word gevind deur die vorige twee getalle in die reeks bymekaar te tel.\\
Die $4^{th}$ word gevind deur die vorige twee getalle in die reeks bymekaar te tel $1+1=2$.\\
Die $9^{th}$  word gevind deur die twee
getalle voor die 2 in die reeks bymekaar te tel $8+13=21$.\\
Die volgende getal in die reeks sal $21+34=55$ wees.
Kan jy die volgende paar getalle vind?

\setcounter{subfigure}{0}
\begin{figure}[H] 
\textnormal{Khan Academy video on Number Patterns - 1}\vspace{.1in} \nopagebreak
\raisebox{-5 pt}{ \includegraphics[width=0.5cm]{col11306.imgs/summary_www.png}} { (Video:  MG10023 )}
\end{figure}       

\begin{wex}{Studeertafels}{Gestel jy en $3$ vriende besluit om te
studeer vir wiskunde, en dat julle om ’n vierkantige tafel sit. ’n Paar
minute later sluit $2$ ander vriende by julle aan en hulle wil kom sit. Om
sitplek te kry vir hulle, besluit julle om ’n tafel te skuif en dit langs julle
tafel te sit sodat daar genoeg sitplek is vir die $6$ van julle. Daarna besluit
nog $2$ van jou vriende om by julle aan te sluit en julle skuif ’n derde tafel
sodat daar genoeg plek is vir $8$ van julle. \\
\\
Ondersoek hoe die aantal mense om die tafels verband hou met die
aantal tafels. Is there a pattern?
\begin{figure}[H]
\begin{center}
\begin{pspicture}(0,0)(9.6,1.8)
%\psgrid[gridcolor=lightgray]
\psframe(0.4,0.4)(1.4,1.4)
\psframe(0,0.6)(0.2,1.2)
\psframe(0.6,0)(1.2,0.2)
\psframe(0.6,1.6)(1.2,1.8)
\psframe(1.6,0.6)(1.8,1.2)
\rput(2.4,0){\psframe(0.4,0.4)(1.4,1.4)
\psframe(0,0.6)(0.2,1.2)
\psframe(0.6,0)(1.2,0.2)
\psframe(0.6,1.6)(1.2,1.8)}
\rput(3.4,0){\psframe(0.4,0.4)(1.4,1.4)
\psframe(0.6,0)(1.2,0.2)
\psframe(0.6,1.6)(1.2,1.8)
\psframe(1.6,0.6)(1.8,1.2)}
\rput(5.8,0){\psframe(0.4,0.4)(1.4,1.4)
\psframe(0,0.6)(0.2,1.2)
\psframe(0.6,0)(1.2,0.2)
\psframe(0.6,1.6)(1.2,1.8)}
\rput(6.8,0){\psframe(0.4,0.4)(1.4,1.4)
\psframe(0.6,0)(1.2,0.2)
\psframe(0.6,1.6)(1.2,1.8)}
\rput(7.8,0){\psframe(0.4,0.4)(1.4,1.4)
\psframe(0.6,0)(1.2,0.2)
\psframe(0.6,1.6)(1.2,1.8)
\psframe(1.6,0.6)(1.8,1.2)}
\end{pspicture}
\caption{Twee ekstra mense kan sit vir elke tafel wat hulle bysit.}
\label{fig:mp:s:arithmetictables}
\end{center}
\end{figure}
}
{
\westep{Make a table to see if a pattern forms}
\begin{center}
\begin{tabular}{|c|l|}\hline
\hline \textbf{Aantal tafels}, $n$ & \textbf{Aantal mense wat sitplek het}\\\hline
\hline 1 & $4 = 4$\\
\hline 2 & $4 + 2 = 6$\\
\hline 3 & $4 + 2 + 2 = 8$\\
\hline 4 & $4 + 2 + 2 + 2 = 10$ \\
\hline \vdots & \qquad \qquad \quad \vdots\\
\hline $n$ & $4 + 2 + 2 + 2 + \ldots + 2 $\\
\hline\hline
\end{tabular}
\end{center}
\westep{Describe the pattern}
Ons kan sien dat met $3$ tafels is daar plek vir $8$ mense, met $4$  tafels
is daar plek vir $10$  mense ens. Ons begin met $4$ mense en voeg
elke keer $2$ mense by.So, vir elke tafel wat bygevoeg word, is daar
sitplek vir nog $2$ mense.}%\end{wex}

\end{wex}

\section{Describing sequences}

 $ \hspace{-5pt}\begin{array}{cccccccccccc}   \includegraphics[width=0.75cm]{col11306.imgs/summary_fullmarks.png} &   \includegraphics[width=0.75cm]{col11306.imgs/summary_video.png} &   \end{array} $ \hspace{2 pt}\raisebox{-5 pt}{} {(section shortcode: MG10024 )} \par 

\setcounter{subfigure}{0}
\begin{figure}[H]
\textnormal{Khan Academy video on Number Patterns}\vspace{.1in} \nopagebreak
% \label{m39362*yt-media3}\label{m39362*yt-video3}
\raisebox{-5 pt}{ \includegraphics[width=0.5cm]{col11306.imgs/summary_www.png}} { (Video:  MG10025 )}
\end{figure}       
 
To describe terms in a number pattern we use the following notation:\\
\\
Die $1^{st}$ term van 'n reeks is $T_{1}$\\
Die $4^{th}$  term van 'n reeks is $T_{4}$\\
Die $10^{th}$  term van 'n reeks is $T_{10}$\\
\\
Die  ${n}^{th}$  van ’n reeks word geskryf as ${T}_{n}$. \\
\\n Reeks hoef nie ’n patroon te volg nie, maar wanneer dit wel ’n patroon het, kan ons
dit gewoonlik as ’n formule skryf om die ${n}^{th}$ - term,${T}_{n}$ te bereken. 
Byvoorveeld,
     
\begin{equation*}
1;3;5;7;9;\ldots
\end{equation}
The ${n}^{th}$ term is given by the general formula
\begin{equation*}
  \begin{array}{cc}\hfill {T}_{n}={2n-1}\end{array}
\end{equation}
You can check this by substituting values into the formula:\
\begin{equation*}
    \begin{array}{ccc}\hfill {T}_{1}& =& 2(1)-1 = 1 \hfill \\ 
    \hfill {T}_{2}& =& 2(2)-1 = 3\hfill \\
    \hfill {T}_{3}& =& 2(3)-1 = 5\hfill \\
    \hfill {T}_{4}& =& 2(4)-1 = 7\hfill \\
    \hfill {T}_{5}& =& 2(5)-1 = 9\hfill \\
\end{array}
\end{equation}

If we find the relationship between the position of a term and its value, we can describe the pattern and find any term in the sequence.

\subsection*{Konstante verskil}
Ons kan ook ’n konstante verskil tussen die terme bepaal vir sekere patrone. 
\\This is called a \textit{Konstante verskil}.

\Definition{Konstante verskil} {Die konstante verskil is die verskil tussen opeenvolgende terme en word aagedui met die
letter $d$. } 

Byvoorbeeld, beskou die reeks $10;7;4;1;\ldots$\\
 Om die gemeenskaplike verskil te vind, trek ons die betrokke term af van die volgende term:

\begin{equation*}
    \begin{array}{ccl}d &=& T_{2} - T_{1}\\
      & =& 7-10\\& =&-3\end{array}
\end{equation}

Let us check another two terms:
\begin{equation*}
    \begin{array}{ccl}d &=& T_{4} - T_{3}\\
    & =& 1-4\\& =&-3\end{array}
\end{equation}
Ons sien dat $d$ konstant is.

\Tip{$d=T_{2}-T_{1}$, nie $T_{1} -T_{2}$}
       
\begin{wex}{Studeertafel voortgesit...}{Soos voorheen, studeer jy
en $3$ vriende wiskunde, en julle sit rondom ’n vierkantige tafel. ’n Paar
minute later besluit $2$ ander vriende om by julle aan te sluit en wil kom sit
en julle sit ’n ekstra tafel by sodat al $6$ van julle kan sit. Weereens besluit
nog $2$ van jou vriende om by julle aan te sluit en julle skuif ’n derde tafel
sodat daar genoeg plek is vir $8$ van julle soos in die prentjie.\\
\\
Vind ’n wiskundige uitdrukking vir die getal mense wat om $n$ tafels kan sit. \\
Gebruik dan die algemene formule om te bepaal hoeveel mense om
$12$ tafels kan sit en\\
hoeveel tafels is nodig sodat $20$ mense kan sit?\\}{

\begin{figure}[H]
\begin{center}
\begin{pspicture}(0,0)(9.6,1.8)
%\psgrid[gridcolor=lightgray]
\psframe(0.4,0.4)(1.4,1.4)
\psframe(0,0.6)(0.2,1.2)
\psframe(0.6,0)(1.2,0.2)
\psframe(0.6,1.6)(1.2,1.8)
\psframe(1.6,0.6)(1.8,1.2)
\rput(2.4,0){\psframe(0.4,0.4)(1.4,1.4)
\psframe(0,0.6)(0.2,1.2)
\psframe(0.6,0)(1.2,0.2)
\psframe(0.6,1.6)(1.2,1.8)}
\rput(3.4,0){\psframe(0.4,0.4)(1.4,1.4)
\psframe(0.6,0)(1.2,0.2)
\psframe(0.6,1.6)(1.2,1.8)
\psframe(1.6,0.6)(1.8,1.2)}
\rput(5.8,0){\psframe(0.4,0.4)(1.4,1.4)
\psframe(0,0.6)(0.2,1.2)
\psframe(0.6,0)(1.2,0.2)
\psframe(0.6,1.6)(1.2,1.8)}
\rput(6.8,0){\psframe(0.4,0.4)(1.4,1.4)
\psframe(0.6,0)(1.2,0.2)
\psframe(0.6,1.6)(1.2,1.8)}
\rput(7.8,0){\psframe(0.4,0.4)(1.4,1.4)
\psframe(0.6,0)(1.2,0.2)
\psframe(0.6,1.6)(1.2,1.8)
\psframe(1.6,0.6)(1.8,1.2)}
\end{pspicture}
\caption{ Twee ekstra mense kan sit vir elke tafel wat hulle bysit.}
\label{fig:mp:s:arithmetictables2}
\end{center}
\end{figure}

\westep{Make a table to see the pattern}
\begin{center}
\begin{tabular}{|c|l|c|}\hline
\hline \textbf{Aantal tafels}, $n$ & \textbf{Aantal mense wat kan sit} & \textbf{Formule}\\\hline
\hline 1 & $4 = 4$ & $= 4 + 2 (0)$ \\
\hline 2 & $4 + 2 = 6$ & $= 4 + 2 t(1)$ \\
\hline 3 & $4 + 2 + 2 = 8$ & $= 4 + 2 (2)$ \\
\hline 4 & $4 + 2 + 2 + 2 = 10$ & $= 4 + 2(3)$ \\
\hline \vdots & \qquad \qquad \quad \vdots & \vdots \\
\hline $n$ & $4 + 2 + 2 + 2 + \ldots + 2 $ & \: \: \: $= 4 + 2 (n-1)$\\
\hline\hline
\end{tabular}
\end{center}

\westep{Describe the pattern}
Die aantal mense wat rondom $n$ tafels kan sit, is:
\nequ{T_n=4 + 2(n-1)}

\westep{Calculate the $12^{th}$ term}
Deur te kyk na die voorbeeld van die vorige gedeelte, bereken ho-
eveel mense kan rondom $12$ tafels sit. Ons soek vir $T_{12}$, dit is, waar $n = 12$:
\begin{eqnarray*}

T_{12} &=& 4 + 2  (12 - 1) \\
&=& 4 + 2(11) \\
&=& 4 + 22 \\
&=& 26
\end{eqnarray*}
Daarom kan $26$ mense kan rondom  $12$ tafels sit .
\westep{Calculate the number of tables needed to seat $20$ people. In other words, find $n$ if $T_n=20$}
\begin{eqnarray*}
T_n &=& 4 + 2  (n - 1) \\
20 &=& 4 + 2  (n - 1) \\
20 - 4 &=& 2  (n - 1) \\
\frac{16}{2} &=& n - 1 \\
8 + 1 &=& n \\
n &=& 9
\end{eqnarray*}
Daarom is $9$ tafels is nodig sodat $20$ mense kan sit.

}
\end{wex}

Dit is ook belangrik om te let op die verskil tussen $n$ en ${T}_{n}$. $n$ kan gesien word as ’n plekhouer, terwyl ${T}_{n}$ die
waarde is by die plek wat "gehou" word deur $n$. Soos in ons "Studeertafel" voorbeeld, kan  $4$ mense rondom die
eerste tafel (Tabel 1) sit. Dus, by plek $n=1$,  is die waarde van ${T}_{1}=4$ ensovoorts:\par 
    
\begin{center}
\begin{tabular}{|c|c|c|c|c|c|}
\hline $n$ & 1 & 2 & 3 & 4 & \ldots \\
\hline $T_n$ & 4 & 6 & 8 & 10 & \ldots \\
\hline
\end{tabular}
\end{center}


\begin{exercises}{Finding the pattern}
{ 

\begin{enumerate}[noitemsep, label=\textbf{\arabic*}. ] 
\item Write down the next three terms in each of the following sequences:

  \begin{enumerate}[noitemsep, label=\textbf{\alph*}. ] 
  \item $5;~15;~25\ldots$
  \item $-8;-3;~ 2 \ldots$
  \item $30;~27;~24\ldots$
  \end{enumerate}
 \item Hieronder is die algemene formules gegee vir ’n paar reekse. Bereken die terme wat weggelaat is.
  \begin{enumerate}[noitemsep, label=\textbf{\alph*}. ] 
  \item $0;3;\ldots;~15;~24$\hspace{2.2cm}$T_{n}={n}^{2}-1$
  \item $3;2;~1;~0;\ldots;-2$\hspace{2cm}$T_{n}=-n+4$
  \item $-11;\ldots;-7;\ldots;-3$\hspace{1.5cm}$T_{n}=-13+2n$
  \end{enumerate}
     
\item Vind die algemene formule vir die volgende reekse en vind dan ${T}_{10}$, ${T}_{50}$ en ${T}_{100}$
  \begin{enumerate}[noitemsep, label=\textbf{\alph*}. ] 
  \item $2;~5;~8;~11;~14;\ldots$
  \item $0;~4;~8;~12;~16;\ldots$
  \item $2;-1;-4;-7;-10;\ldots$
  \end{enumerate}
\end{enumerate}
}%\End of exercise



\section{Patrone en Bewerings}

\setcounter{subfigure}{0}
\begin{figure}[H] 
\textnormal{Khan Academy video on Number Patterns - 2}\vspace{.1in} k
\raisebox{-5 pt}{ \includegraphics[width=0.5cm]{col11306.imgs/summary_www.png}} { (Video:  MG10026 )}
\end{figure} 

In wiskunde is ’n bewering ’n wiskundige stelling wat lyk of dit waar is, maar wat nog nie formeel as waar bewys
is nie. 
A conjecture can be thought of as the mathematicians way of saying ``I believe that this is true, but I have no proof''.
’n Bewering kan gesien word as ’n intelligente raaiskoot of idee wat moontlik ’n patroon kan wees. \\ 
\\ Byvoorbeeld: Maak ’n bewering oor die getal wat sal volg, gebaseer op die patroon $2;~6;~11;~17;\ldots$\\
Die getalle vermeerder met $4$, dan $5$, en dan $6$.\\
Bewering: Die volgende getal sal vermeerder met $7$, So ons verwag dat die volgende getal $17+7=24$ sal wees.\\

\begin{wex}{Adding even and odd numbers }{
\begin{enumerate}[noitemsep, label=\textbf{\arabic*}. ] 
\item Investigate the type of number you get if you find the sum of an odd number and an even number.
\item Express your answer in words as a conjecture.
\item Use algebra to prove this conjecture.
\end{enumerate}
}
{   
\westep{First try some examples}  
  \begin{equation*}
    \begin{array}{ccc}\hfill 23 + 12 &=& 35\\ 148 + 31 &=& 179\\ 11 +200 &=& 211 \end{array}
   \end{equation}

\westep{Make a conjecture}  
        The sum of any odd number and any even number is always odd.  \\
\westep{Express algebraically}
Express the even number as $2x$.\\
Express the odd number as $2y-1$.\\

\begin{equation*}
    \begin{array}{ccc}\hfill \mbox{Sum} &=& 2x + (2y-1)\\  &=& 2x + 2y - 1\\ &=& (2x + 2y) - 1 \\ &=& 2(x+y) -1 \end{array}
\end{equation}
From this we can see that $2(x+y)$ is an even number. So then $2(x+y)-1$ is an odd number.\\
So our conjecture is true.    
}
\end{wex}



\begin{wex}{Multiplying a two-digit number by $11$ }{
\begin{enumerate}[noitemsep, label=\textbf{\arabic*}. ] 
\item Consider the following examples of multiplying any two-digit number by $11$. What pattern can you see?
\begin{equation*}
    \begin{array}{ccc}\hfill 11 \times 42 &=& 462\\ 11 \times 71 &=& 781\\ 11 \times 45 &=& 495\end{array}
\end{equation}

\item Express your answer as a conjecture.
\item Can you find an example that disproves your initial conjecture? If so, reconsider your conjecture.
\item Use algebra to prove your conjecture.
\end{enumerate}
}
{   
\westep{Finding the pattern}
We notice in the answer that the middle digit is the sum of the two digits in the original two-digit number.  \\

\westep{Make a conjecture}
The middle digit of the product is the sum of the two digits of the original number that is multipled by $11$.\\
   
\westep{Reconsidering the conjecture}
Yes. We notice that 
\begin{equation*}
    \begin{array}{ccc}\hfill 11 \times 67 &=& 737\\ 11 \times 56 &=& 638\end{array}
\end{equation}
Therefore our conjecture only holds true if the sum of the two digits is less than $10$.\\

\westep{beskryf die bewering algebraïes}
Any two-digit number can be written as $10a+b$. Byvoorbeeld $34= 10(3)+4$ \\
Any three-digit number can be written as $100a+10b+c$. Byvoorbeeld $582=100(5)+ 10(8) +2$\\
\begin{equation*}
    \begin{array}{ccl}\hfill 11 \times (10x+y) &=& 110x + 11y\\ &=&(100x + 10x) + 10y + y\\&=& 100x + (10x+ 10y) + y\\&=& 100x + 10(x+y) + y\end{array}
\end{equation} 

From this equation we can see that the middle digit of the three-digit number is equal to the sum of the two digits $x$ and $y$
}
\end{wex}



\section{Opsomming}
\begin{itemize}[noitemsep]
\item Daar is ’n hele paar spesiale reekse van getalle: 
    \begin{itemize}[noitemsep]
    \item Driehoeksgetalle  $1;~3;~6;~10;~15;~21;~28;~36;~45;~...$\
    \item Vierkantsgetalle $1;~4;~9;~16;~25;~36;~49;~64;~81;\ldots$
    \item Derdemagsgetalle $1;~8;~27;~64;~125;~216;~343;~512;~729;\ldots$
    \item Fibonacci Getalle  $0;~1;~1;~2;~3;~5;~8;~13;~21;~34;\ldots$\
    \end{itemize}
\item Die ${n}^{th}$ term word geskryf as ${T}_{n}$
\item Die konstante verskil ($d$) is die verskil tussen opeenvolgende terme
\item We can work out a general formula for each number pattern and use it to determine any term in the pattern.
\item A conjecture is something that you assume to be true but have not yet proved.
\end{itemize}
            

\begin{eocexercises}{Hoofstukoefeninge}
\begin{enumerate}[noitemsep, label=\textbf{\arabic*}. ] 
\item Vind die $6^{th}$ term vir die volgende reekse:
    \begin{enumerate}[noitemsep, label=\textbf{\alph*}. ] 
    \item $4;~13;~22;~31;\ldots$
    \item $~5;~2;~-1;~-4;\ldots$
    \item $7,4;~ 9,7; ~12; ~14,3; \ldots$
    \end{enumerate}


\item Vind die algemene term vir die volgende reekse:
    \begin{enumerate}[noitemsep, label=\textbf{\alph*}. ] 
    \item$3;~7;~11;~15;\ldots$
    \item $-2;~1;~4;~7;\ldots$
    \item $11;~15;~19;~23;\ldots$
    \item $\dfrac{1}{3};~ \dfrac{2}{3};~ 1; ~1\dfrac{1}{3};\ldots$
    \end{enumerate}

\item Die sitplekke in ’n gedeelte van ’n sportstadion kan so gerangskik word dat die eerste ry $15$ sitplekke het,
die tweede ry $19$ sitplekke, die derde ry $23$ sitplekke, ens. Bereken hoeveel sitpleke is daar in ry $25$.
\item ’n Enkele vierkant kan gemaak word van $4$ vuurhoutjies. Om twee vierkante langs mekaar te maak het jy $7$ vuurhoutjies nodig, om drie vierkante langs mekaar in ’n ry te maak het jy $10$ vuurhoutjies nodig. Bepaal:
    \begin{enumerate}[noitemsep, label=\textbf{\alph*}. ] 
    \item die eerste term
    \item die konstante verskil
    \item die algemene formule
    \item hoeveel vuurhoutjies benodig word om $25$ vierkante langs mekaar te maak
    \end{enumerate}

\setcounter{subfigure}{0}
\begin{figure}[H] 
\begin{center}
\begin{pspicture}(0,0)(8,2)
%\psgrid[gridcolor=gray]
\def\match{\psline(0,0)(2,0)\psellipse*(1.8,0)(0.2,0.1)}
\rput(0,0){\match}
\rput{90}(2,0){\match}
\rput{180}(2,2){\match}
\rput{270}(0,2){\match}
\rput(2,0){\rput(0,0){\match}
\rput{90}(2,0){\match}
\rput{180}(2,2){\match}}
\rput(4,0){\rput(0,0){\match}
\rput{90}(2,0){\match}
\rput{180}(2,2){\match}}
\rput(6,0){\rput(0,0){\match}
\rput{90}(2,0){\match}
\rput{180}(2,2){\match}}
\end{pspicture}

\vspace{2pt}
\vspace{.1in}
\end{center}
\end{figure}       

\item Jy wil begin om geld te spaar, maar omdat jy dit nog nooit gedoen het nie, besluit jy om stadig te begin.
Aan die einde van die eerste week sit jy R$~5$ in jou bankrekening, aan die einde van die tweede week R$~10$ en aan die einde van die derde week R$~15$. Na hoeveel weke sit jy
 R$~50$ in jou bankrekening?

\item ’n Horisontale lyn kruis ’n tou op vier punte en deel die tou op in $4$ dele, soos hieronder gewys word.
\setcounter{subfigure}{0}
\begin{figure}[H] 
\begin{center}

\begin{pspicture}(-1,-2)(6,2)
%\psgrid[gridcolor=gray]
\psplot[xunit=0.00556]{90}{810}{x sin}
\psline[linestyle=dashed](-1,0)(6,0)
\psdots[dotsize=5pt](1,0)(2,0)(3,0)(4,0)
\rput(0.5,1.5){\psframebox{1}}
\rput(1.5,-1.5){\psframebox{2}}
\rput(2.5,1.5){\psframebox{3}}
\rput(3.5,-1.5){\psframebox{4}}
\rput(4.5,1.5){\psframebox{5}}
\end{pspicture}
\end{center}
\end{figure}  
     
As die tou $19$ keer gekruis word deur ewewydige lyne en elke lyn kruis die tou $4$ keer op verskillende
plekke, bereken in hoeveel dele die tou opgedeel word.
.
 
\item Consider the following pattern: 
  \begin{equation*}
    \begin{array}{ccl}\hfill 9+16&=& 25\\ 9+28 &=& 37\\9+43&=& 52\end{array}
  \end{equation} 

  \begin{enumerate}[noitemsep, label=\textbf{\alph*}. ] 
  \item What pattern do you see?
  \item Make a conjecture and express it in words.
  \item Generalise your conjecture algebraically.
  \item Prove that your conjecture is true.
  \end{enumerate}

\end{enumerate}


\raisebox{-5 pt}{\includegraphics[width=0.5cm]{col11306.imgs/summary_www.png}} Find the answers with the shortcodes:

\begin{tabular}[h]{cccccc}
 (1.) lcq  &  (2.) lcl  &  (3.) lci  &  (4.) lOe  &  (5.) lcO  &  (6.) lcc & (7.)

\end{tabular}

\end{eocexercises}
